\documentclass{article}
\usepackage[margin=1in]{geometry}
\usepackage{amsmath,amssymb,amsthm}
\usepackage{booktabs,longtable}
\usepackage{hyperref}

\newtheorem{definition}{Definition}
\newtheorem{theorem}[definition]{Theorem}
\newtheorem{proposition}[definition]{Proposition}
\newtheorem{lemma}[definition]{Lemma}
\newtheorem{remark}[definition]{Remark}

\newcommand{\source}[1]{\par\noindent\textbf{Source.} #1}
\newcommand{\lean}[1]{\par\noindent\textbf{Lean target.} \texttt{\detokenize{#1}}}
\newcommand{\status}[1]{\par\noindent\textbf{Status.} \texttt{\detokenize{#1}}}
\newcommand{\depends}[1]{\par\noindent\textbf{Depends on.} #1}
\newcommand{\cA}{\mathcal A}
\newcommand{\cB}{\mathcal B}
\newcommand{\cD}{\mathcal D}
\newcommand{\cH}{\mathcal H}
\newcommand{\cM}{\mathcal M}
\newcommand{\cP}{\mathcal P}
\newcommand{\cR}{\mathcal R}

\begin{document}

\section*{Blueprint 012: Infinite Products and Tensor Type Calculus}

\subsection*{Scope}

This note promotes the fifth remaining blueprint block from
\path{agents/2026-06-22-remaining-blueprint-suite-goal.md}.  It supports
\[
  \texttt{MFH-001},\quad
  \texttt{ITP-001},\quad
  \texttt{ITP-PURE-001},\quad
  \texttt{ITP-TRACE-001},\quad
  \texttt{TENSOR-001}.
\]
It supplies the infinite product and standard-factor type calculus consumed by
the final stabilizer type-classification theorem.  It does not introduce Lean
files or Aristotle packets.

\subsection*{Source Policy for This Block}

Takesaki III gives the primary infinite tensor product construction relative to
normal product states and reference vectors.  Takesaki I gives the project-local
type vocabulary, finite matrix and \(\cB(\cH)\) model factors, trace
interfaces, and finite tensor-product representation bicommutants.  The
specific statements that the countable tracial \(M_2\) product is the
hyperfinite \(II_1\) factor \(\cR\), and that
\(\cR\bar\otimes \cB(\cH)\) is type \(II_\infty\), were not located in the
checked Takesaki locators.  They are recorded below either as precise
Blackadar-backed theorem targets with proof routes, or as explicitly named
conditional hypotheses that later theorem statements may assume if the local
Lean development has not yet formalized those targets.

\begin{definition}[InfiniteTensorProductVNStates, Takesaki III Definition XIV.1.6, PDF p.~107]
Given a sequence \(\{\cM_n\}\) of von Neumann algebras and normal states
\(\omega_n\) of \(\cM_n\), define the product state
\(\omega=\bigotimes_{n=1}^{\infty}\omega_n\) on the \(C^*\)-algebraic
infinite tensor product by
\[
  \omega(x_1\otimes x_2\otimes\cdots\otimes x_n\otimes 1\otimes 1\otimes
  \cdots)
  =
  \omega_1(x_1)\omega_2(x_2)\cdots\omega_n(x_n).
\]
Let \(\{\pi_\omega,\cH_\omega,\xi_\omega\}\) be the cyclic representation of
this product state.  The von Neumann algebra
\[
  \pi_\omega\!\left(\bigotimes_{n=1}^{\infty}\cM_n\right)''
  \quad\hbox{on }\cH_\omega
\]
is called the infinite tensor product of \(\{\cM_n\}\) relative to the states
\(\{\omega_n\}\).
\end{definition}

\begin{proposition}[TensorProductRepresentationBicommutant, Takesaki I Proposition IV.4.13, PDF p.~231]
Let \(A_1\) and \(A_2\) be \(C^*\)-algebras, and let
\(\{\pi_1,\mathcal H_1\}\) and \(\{\pi_2,\mathcal H_2\}\) be their respective
representations.  Then
\[
  \{(\pi_1\otimes\pi_2)(A_1\otimes_{\min}A_2)\}''
  =
  \pi_1(A_1)''\bar\otimes\pi_2(A_2)'',
\]
where the latter means the tensor product of the von Neumann algebras
\(\pi_1(A_1)''\) and \(\pi_2(A_2)''\).  In particular, if \(\pi_1\) and
\(\pi_2\) are both irreducible, then so is \(\pi_1\otimes\pi_2\).
\end{proposition}

\begin{proof}
The represented minimal tensor product is the uniform closure of the algebraic
tensor product \(\pi_1(A_1)\otimes\pi_2(A_2)\), so it sits inside the von
Neumann tensor product of the bicommutants.  Conversely, approximate each
operator in \(\pi_1(A_1)''\) and \(\pi_2(A_2)''\) \(\sigma\)-strongly by
represented elements and use the \(\sigma\)-strong continuity of tensoring by
the identity to put the generators \(x_1\otimes 1\) and \(1\otimes x_2\) in
the bicommutant.  The irreducible case follows because the bicommutants are
full operator algebras.
\end{proof}

\begin{definition}[ModelFactorConditionalHypotheses, ID: MFH-001, Blackadar III.3.1.2--III.3.1.6 and III.3.4.1--III.3.4.2, PDF pp.~301--303, 317]
If a later Lean extraction reaches the final type theorem before the
Blackadar-backed model-factor targets in this note have been formalized, the
only permitted high-level assumptions are the following named propositions:
\begin{description}
\item[\(\mathrm{MFH}_{\mathrm{factor}}\).]
For every finite or countably infinite family
\((\cM_i,\omega_i)_{i\in I}\) of factors with faithful normal reference
states, the represented product
\(\overline{\bigotimes}_{i\in I}(\cM_i,\omega_i)\) is a factor, and if all
\(\omega_i\) are traces then the product GNS vector state is a faithful normal
trace on that represented product.
\item[\(\mathrm{MFH}_{\mathrm{tracial}\,M_2}\).]
For \(R\) countably infinite, the product
\[
  \overline{\bigotimes}_{i\in R}(M_2(\mathbb C),\operatorname{tr}_2)
\]
is the hyperfinite type \(II_1\) factor \(\cR\), with faithful normal
tracial vector state induced by \(\bigotimes_R\operatorname{tr}_2\).
\item[\(\mathrm{MFH}_{\cR\otimes\cB}\).]
For a separable infinite-dimensional Hilbert space \(\cH\),
\(\cR\bar\otimes\cB(\cH)\) is a type \(II_\infty\) factor and admits the
faithful semifinite normal trace obtained from the finite trace on \(\cR\)
and the standard semifinite trace on \(\cB(\cH)\).
\end{description}
These assumptions are model-factor statements only.  They do not identify a
stabilizer region algebra with any model factor; that identification must come
separately from \texttt{PROD-001}, \texttt{GNS-003}, and \texttt{VN-003}.
\source{Blackadar, \emph{Operator Algebras}, Sections III.3.1.2--III.3.1.6
and III.3.4.1--III.3.4.2, local PDF \path{../sources/blackadar.pdf}, PDF
pp.~301--303 and 317.}
\depends{Blackadar model-factor source statements; consumed by
\texttt{ITP-001}, \texttt{ITP-TRACE-001}, and \texttt{TENSOR-001} when those
proofs are deferred.}
\lean{Formalization.OperatorAlgebra.TensorTypeCalculus.ModelFactorHypotheses}
\status{named_conditional_hypotheses_if_model_factor_proofs_are_deferred}
\end{definition}

\begin{proof}[Well-formedness]
Each named hypothesis fixes its ambient von Neumann algebra, reference state
or trace, and target type predicate before it is used.  The hypotheses are
closed operator-algebra assertions about standard product models, so a later
formal theorem can quantify them explicitly without hiding stabilizer-specific
content inside a source citation.
\end{proof}

\subsection*{Product-State Infinite Tensor Products}

\begin{definition}[ProductStateVN, ID: ITP-001, Takesaki III Definition XIV.1.6, PDF p.~107]
Let \(I\) be finite or countably infinite.  For each \(i\in I\), let
\(\cM_i\) be a von Neumann algebra with normal state \(\omega_i\).  The
project infinite product package consists of:
\begin{itemize}
\item the C*-algebraic infinite tensor product when \(I\) is infinite and the
usual finite tensor product when \(I\) is finite;
\item the product state \(\bigotimes_{i\in I}\omega_i\);
\item the GNS Hilbert space and cyclic representation of that product state;
\item the represented von Neumann algebra
\[
  \overline{\bigotimes}_{i\in I}(\cM_i,\omega_i)
  :=
  \pi_{\otimes_i\omega_i}\!\left(\bigotimes_i \cM_i\right)'' .
\]
\end{itemize}
\source{Takesaki III, Definition XIV.1.6, PDF p.~107, for the construction
relative to normal states; Blackadar III.3.1.2, PDF pp.~301--302, for the
product GNS and trace-preservation theorem route used in the stabilizer
model.}
\depends{\texttt{GNS-003}, \texttt{VN-001}, \texttt{VN-003}.}
\lean{Formalization.OperatorAlgebra.InfiniteTensorProduct.productStateVN}
\status{blueprint_note_complete}
\end{definition}

\begin{proposition}[InfiniteTensorProductFactorProduct, supporting ID: ITP-001, Blackadar III.3.1.3, PDF p.~302]
If each \((\cM_i,\omega_i)\) in the product package is a factor with faithful
normal reference state, then
\[
  \overline{\bigotimes}_{i\in I}(\cM_i,\omega_i)
\]
is a factor.
\source{Blackadar, Proposition III.3.1.3, PDF p.~302.  Takesaki supplies the
underlying relative-state product construction in Definition XIV.1.6, PDF
p.~107.}
\depends{\texttt{ITP-001}, \texttt{VN-002}.}
\lean{Formalization.OperatorAlgebra.InfiniteTensorProduct.factor_product}
\status{blueprint_note_complete}
\end{proposition}

\begin{proof}
For finite \(I\), apply the finite tensor-product bicommutant theorem and the
fact that a finite tensor product of factors has scalar center.  For countable
\(I\), let \(E_F\) denote the normal product-state conditional expectation
onto the finite subproduct over \(F\Subset I\).  If \(z\) is central in the
represented infinite product, then \(E_F(z)\) is central in the finite
subproduct and hence scalar.  The finite subproducts are increasing and
\(\sigma\)-weakly dense in the product representation; normality of the
expectations and density force \(z\) to agree with the scalar limit of
\(E_F(z)\).  Thus the center is scalar.  Until this conditional-expectation
argument is formalized locally, later theorem statements may assume exactly
the \(\mathrm{MFH}_{\mathrm{factor}}\) clause of \texttt{MFH-001}.
\end{proof}

\subsection*{Pure-Mode Product Factor}

\begin{proposition}[PureModeProductTypeI, ID: ITP-PURE-001, Takesaki III Definition XIV.1.6, PDF p.~107, Blackadar III.3.1.2, PDF pp.~301--302]
Let \(\omega_{Z,+}\) be the pure \(+1\)-eigenstate of \(Z\) on
\(M_2(\mathbb C)\), and let
\[
  \cP_P =
  \overline{\bigotimes}_{i\in P}\bigl(M_2(\mathbb C),\omega_{Z,+}\bigr)
\]
for a finite or countably infinite set \(P\).  If \(P\) is finite, then
\[
  \cP_P \cong M_{2^{|P|}}(\mathbb C),
\]
and if \(P\) is countably infinite, then
\[
  \cP_P \cong \cB(\ell^2).
\]
Thus \(\cP_P\) is finite type \(I_{2^{|P|}}\) for finite \(P\), and type
\(I_\infty\) for countably infinite \(P\).
\source{Takesaki III, Definition XIV.1.6, PDF p.~107, for the relative-state
infinite tensor product; Blackadar III.3.1.2, PDF pp.~301--302, for the
commutant statement in finite-dimensional tensor products; Takesaki I,
Proposition V.1.23, Corollary V.1.28, and Definition V.1.33, PDF pp.~318--322,
for finite matrix and \(\cB(\cH)\) type I model factors.}
\depends{\texttt{ITP-001}, \texttt{MATRIX-001}.}
\lean{Formalization.OperatorAlgebra.PureModeProduct.typeI}
\status{blueprint_note_complete}
\end{proposition}

\begin{proof}
For a single site, the GNS representation of a pure vector state on
\(M_2(\mathbb C)\) is the defining irreducible representation on
\(\mathbb C^2\); its commutant is scalar.  Finite products give the defining
representation of \(M_{2^{|P|}}(\mathbb C)\).  In the countable case, use the
infinite tensor product Hilbert space relative to the reference \(+1\)
eigenvectors.  The commutant of the product algebra is the product of the
scalar one-site commutants, hence scalar, so the bicommutant is
\(\cB(\cH_P)\).  The incomplete tensor product Hilbert space for countably
many two-dimensional sites and a reference vector is separable infinite
dimensional, hence isomorphic to \(\ell^2\).  The type conclusion is then
\texttt{MATRIX-001}.
\end{proof}

\subsection*{Tracial Defect Product}

\begin{proposition}[TracialDefectProductHyperfiniteII1, ID: ITP-TRACE-001, Blackadar III.3.1.4 and III.3.4.2, PDF pp.~302, 317]
Let \(\operatorname{tr}_2\) be the normalized trace on \(M_2(\mathbb C)\), and
let
\[
  \cD_R =
  \overline{\bigotimes}_{i\in R}\bigl(M_2(\mathbb C),\operatorname{tr}_2\bigr)
\]
for a finite or countably infinite set \(R\).  If \(R\) is finite, then
\[
  \cD_R \cong M_{2^{|R|}}(\mathbb C)
\]
with its normalized finite trace.  If \(R\) is countably infinite, then
\(\cD_R\) is the hyperfinite type \(II_1\) factor \(\cR\).
\source{Takesaki III, Definition XIV.1.6, PDF p.~107, for the relative-state
infinite tensor product; Blackadar III.3.1.4, PDF p.~302, for the tracial
\(M_2\) product as a \(II_1\) factor; Blackadar III.3.4.1--III.3.4.2, PDF
p.~317, for hyperfiniteness of infinite tensor products of finite-dimensional
von Neumann algebras and hence of \(\cR\).}
\depends{\texttt{ITP-001}, \texttt{TRACE-001}, \texttt{MATRIX-001}.}
\lean{Formalization.OperatorAlgebra.TracialDefectProduct.hyperfiniteII1}
\status{blueprint_note_complete}
\end{proposition}

\begin{proof}
For finite \(R\), the algebra is the finite matrix tensor product, hence
isomorphic to \(M_{2^{|R|}}(\mathbb C)\), and the product state is the
normalized matrix trace.  For countably infinite \(R\), the one-site trace is
faithful and normal, so \(\mathrm{MFH}_{\mathrm{factor}}\) gives factoriality
and a faithful normal tracial vector state on the product.  Blackadar
III.3.1.4 identifies this tracial product as the \(II_1\) factor \(R_1\).
Blackadar III.3.4.1--III.3.4.2 give a directed family of finite-dimensional
subalgebras with \(\sigma\)-weakly dense union, so the factor is hyperfinite.
Until that source-backed identification is formalized locally, later theorem
statements may assume exactly the \(\mathrm{MFH}_{\mathrm{tracial}\,M_2}\)
clause of \texttt{MFH-001}.
\end{proof}

\subsection*{Pure--Defect Tensor Type Calculus}

\begin{theorem}[PureDefectTensorTypeCalculus, ID: TENSOR-001, Takesaki I Proposition IV.4.13, PDF p.~231, Blackadar III.3.1.6, PDF p.~303]
Let \(\cP_P\) and \(\cD_R\) be the pure-mode and tracial-defect model factors
above.  The represented stabilizer product factor is
\[
  \cP_P \bar\otimes \cD_R .
\]
Its type is:
\[
\begin{array}{c|c|c}
|P|<\infty & |R|<\infty & \hbox{finite type } I\\
|P|=\infty & |R|<\infty & \hbox{type } I_\infty\\
|P|<\infty & |R|=\infty & \hbox{hyperfinite type } II_1\\
|P|=\infty & |R|=\infty & \hbox{type } II_\infty .
\end{array}
\]
\source{Takesaki I, Proposition IV.4.13, PDF p.~231, for represented tensor
product bicommutants; Takesaki I, Proposition V.1.23, Corollary V.1.28,
Definition V.1.33, Definition V.2.1, and Theorem V.2.15, PDF pp.~318--337,
for matrix/\(\cB(\cH)\) model factors and finite/semifinite trace interfaces;
Blackadar III.3.1.4 and III.3.4.2, PDF pp.~302, 317, for the hyperfinite
\(II_1\) factor; Blackadar III.3.1.6, PDF p.~303, for
\(\cR\bar\otimes\cB(\cH)\) as type \(II_\infty\).}
\depends{\texttt{ITP-PURE-001}, \texttt{ITP-TRACE-001},
\texttt{TRACE-001}, \texttt{TYPE-DEF-001}, \texttt{TYPE-ALT-001},
\texttt{MATRIX-001}; if the model-factor proofs are deferred,
\texttt{MFH-001}, namely the named hypotheses
\(\mathrm{MFH}_{\mathrm{factor}}\),
\(\mathrm{MFH}_{\mathrm{tracial}\,M_2}\), and
\(\mathrm{MFH}_{\cR\otimes\cB}\).}
\lean{Formalization.OperatorAlgebra.TensorTypeCalculus.pure_defect_type}
\status{blueprint_note_complete}
\end{theorem}

\begin{proof}
Combine the model identifications from \texttt{ITP-PURE-001} and
\texttt{ITP-TRACE-001}.  If \(P\) and \(R\) are finite, both factors are finite
matrix algebras, so the tensor product is again a finite matrix factor and is
finite type I by \texttt{MATRIX-001}.  If \(P\) is countably infinite and
\(R\) is finite, then \(\cP_P\cong\cB(\ell^2)\) and tensoring with a finite
matrix algebra only replaces \(\ell^2\) by an isomorphic separable
infinite-dimensional Hilbert space; \texttt{MATRIX-001} gives type
\(I_\infty\).

If \(P\) is finite and \(R\) is countably infinite, then
\(\cP_P\) is finite matrix and \(\cD_R\cong\cR\) by
\texttt{ITP-TRACE-001}.  The product normalized trace is faithful and normal,
and finite matrix amplification preserves the \(II_1\) factor property and
hyperfiniteness.  If both \(P\) and \(R\) are countably infinite, then
\[
  \cP_P\bar\otimes\cD_R
  \cong
  \cB(\ell^2)\bar\otimes\cR
  \cong
  \cR\bar\otimes\cB(\ell^2).
\]
The trace from \(\cR\) tensored with the standard semifinite trace on
\(\cB(\ell^2)\) is faithful, semifinite, and normal, so
\texttt{TRACE-001} gives semifiniteness; Blackadar III.3.1.6, or exactly
the \(\mathrm{MFH}_{\cR\otimes\cB}\) clause of \texttt{MFH-001} if deferred,
identifies the factor type as \(II_\infty\).
\end{proof}

\subsection*{Lean Extraction Notes}

\begin{longtable}{@{}p{0.18\linewidth}p{0.34\linewidth}p{0.38\linewidth}@{}}
\toprule
Target & Intended file & Notes \\
\midrule
\endhead
\texttt{MFH-001} &
\texttt{Formalization/OperatorAlgebra/TensorTypeCalculus.lean} &
Expose only the three named Blackadar-backed model-factor hypotheses when the
local Lean proofs of those model identifications are deferred. \\
\texttt{ITP-001} &
\texttt{Formalization/OperatorAlgebra/InfiniteTensorProduct.lean} &
Start as a narrowly scoped product-state package for finite/countable products;
do not attempt a general von Neumann algebra inductive-limit theory. \\
\texttt{ITP-PURE-001} &
\texttt{Formalization/OperatorAlgebra/PureModeProduct.lean} &
The countable pure product can be formalized through scalar commutant and
irreducibility, then identified with \(\cB(\ell^2)\). \\
\texttt{ITP-TRACE-001} &
\texttt{Formalization/OperatorAlgebra/TracialDefectProduct.lean} &
Formalize the finite matrix case directly; if hyperfinite \(II_1\)
infrastructure is absent, expose \(\mathrm{MFH}_{\mathrm{tracial}\,M_2}\) as
an explicit theorem hypothesis. \\
\texttt{TENSOR-001} &
\texttt{Formalization/OperatorAlgebra/TensorTypeCalculus.lean} &
Formalize only the four stabilizer model-factor cases before any general tensor
type theorem, with \(\mathrm{MFH}_{\cR\otimes\cB}\) explicit if the
Blackadar model proof is deferred. \\
\bottomrule
\end{longtable}

\begin{thebibliography}{9}
\bibitem{TakesakiI}
M. Takesaki, \emph{Theory of Operator Algebras I}.  Local source:
\path{../sources/takesaki_1.pdf}.

\bibitem{TakesakiIII}
M. Takesaki, \emph{Theory of Operator Algebras III}.  Local source:
\path{../sources/Takesaki_III.pdf}.

\bibitem{Blackadar}
B. Blackadar, \emph{Operator Algebras}.  Local source:
\path{../sources/blackadar.pdf}.
\end{thebibliography}

\end{document}
